\section{Conclusions}
\label{sec:conclusions}

We presented RTXGNN, a self-explaining temporal graph neural network for fraud detection. Its masks gate message passing and are returned as the explanation of every score, and its training objective makes the top-$k$ explanation sufficient and its complement uninformative. The temporal encoding (HRAPE) uses only relative time. In a controlled comparison with 16 baselines over 10 seeds on Elliptic, RTXGNN matches the best graph neural network (F1 \rtxFone, AP \rtxAP) and is ahead of the self-explaining SEFraud on average. Its ranking is precise at the top of the alert queue, which is where analysts spend their effort. Gradient-boosted trees on Elliptic's engineered features remain more accurate (F1 \xgbFone). An equal-budget hyperparameter search for RTXGNN and ten baselines leaves this ranking unchanged. They do not, however, produce explanations over the relations between transactions. RTXGNN's explanations come with every score, without a separate explainer and at a small cost. They are nearly as sufficient as Integrated Gradients and much more stable than GNNExplainer, and on synthetic data they single out the planted laundering transactions. A single transaction is scored and explained end to end in \rtxLatBoneP{} ms (99th percentile) on one CPU core.

The evaluation also shows the limits of the approach. All models fail after the regime change at time step 43, and only retraining on newly labelled data recovers part of the performance. The explanations are weakly necessary: removing the top features rarely changes the prediction, because Elliptic's features are redundant. HRAPE carries little information on Elliptic's snapshots and is unavailable on T-Finance, so its benefit rests mainly on comparisons with Time2Vec and on synthetic evidence. Finally, the explanations support, but do not by themselves establish, the transparency expected by regulation. That requires governed feature definitions, curated reason codes, human review and validation with analysts.

Future work should (i) evaluate the temporal components on real transaction data with continuous timestamps and repeated account activity, (ii) study continual-learning and drift-detection strategies that shorten the recovery after regime changes, (iii) conduct a user study with fraud analysts on the usefulness of the explanations, and (iv) assess the robustness of models and explanations against adversaries who know the explanation mechanism.
